\documentclass[11pt]{article}
\usepackage[letterpaper,margin=1in]{geometry}
\usepackage{graphicx}
\usepackage{booktabs}
\usepackage{array}
\usepackage{amsmath}
\usepackage{xcolor}
\usepackage[hidelinks]{hyperref}
\usepackage[capitalise,nameinlink]{cleveref}
\usepackage{caption}
\captionsetup{font=small,labelfont=bf}
\graphicspath{{../results/figures_nocap/}}

% Every number below carries a "% src:" comment naming the table it comes from, under
% results/tables/ unless stated. Regenerate the tables with analysis/run_all.py.

\newcommand{\pmrep}{\texttt{.pm} replay}

\title{Evaluating Halo-Exchange Stand-Ins for Conjugate Gradient\\and the Statistics Used to Judge Them}
\author{Nicholas Bacon\\University of New Mexico}
\date{October 2026}

\begin{document}
\maketitle

\begin{abstract}
Choosing a halo-exchange strategy for a sparse solver from a cheaper measurement, such as a
communication benchmark or the solver's own exchange run alone, is worthwhile only if the
cheaper measurement ranks the strategies the way the solver does and if the statistics
used to compare strategies can be trusted. However, benchmark run times on shared HPC
systems are rarely normally distributed, and the stand-ins differ from the solver in code
path as well as in what they leave out. In this paper we measure both questions on LLNL's
Dane system, using 40 separate allocations at each of 2, 4, 8, and 16 nodes, six
SuiteSparse matrices, three strategies (plain MPI point-to-point, MPI Advance neighbor
alltoallv, and MPI Advance locality-aware aggregation), and seven ways of running the SpMV
halo exchange, from a replay of Vernier's job-wide histograms up to ABCG's conjugate
gradient (CG) iteration. Normality is rejected in 86\% of configurations on run time and
in 62\% on its logarithm, largely because MPI and STANDARD runs lose whole multiples of
0.160\,s. A paired $t$-test on log time nevertheless keeps close to its nominal 5\% false
positive rate and replicates its verdicts as well as rank-based tests do, while a $t$-test
on raw time does not, and with stalls present no test detects a true 10\% difference
reliably with 40 allocations. ABCG's standalone SpMV exchange reaches the same verdict as
CG on 62 of 72 strategy pairs, and every disagreement but two is between MPI and
STANDARD; the matrix replay reaches CG's verdict on 49, because its locality-aware path is
1.49 to 1.99 times slower than ABCG's on the same messages.
\end{abstract}

\section{Introduction}

No single halo-exchange strategy is fastest for every sparse workload. Collom, Li, and
Bienz showed that the best of direct point-to-point exchange and locality-aware
aggregation changes across the levels of one algebraic multigrid
hierarchy~\cite{collom2023}, and the research plan this work belongs to proposes choosing
a strategy once per program launch from cheaply measured pattern
features~\cite{researchplan}. Such a choice needs two things: a stand-in that runs faster
than the application but ranks the strategies the same way, and a statistical procedure
that says when one strategy is faster than another on a machine whose run times vary from
allocation to allocation. A companion report measured how much ABCG's CG iteration slows
down when each of four stand-ins makes the choice instead of CG
itself~\cite{ladderreport}.

However, that report scored stand-ins by comparing medians and did not ask whether the
statistics behind the comparison were sound, and it left open why the stand-ins disagree
with CG where they do. Both questions matter beyond this study. The benchmarking
literature warns against assuming normality without checking it and recommends
nonparametric summaries~\cite{hoefler2015scientific,maricq2018taming}, but a diagnostic
check does not say how much a normal-theory test actually errs on a given workload, which
decides whether its simpler analysis is acceptable. Likewise, a communication benchmark
used as a stand-in for an application runs a different code path and leaves out the
application's computation, and which of those differences changes a strategy comparison
is rarely measured.

This paper evaluates, on the same 40-allocation campaign, which statistical tests give
reliable verdicts on halo-exchange strategy comparisons, and how well each stand-in
reproduces the verdicts of the run above it. By \emph{stand-in} we mean a cheaper
measurement that exchanges a version of the SpMV halo of the same matrix at the same node
count; we do not consider analytical models or simulators. This paper makes the following
contributions:
\begin{itemize}
\item A comparison, line by line, of what the replay benchmark, ABCG's SpMV exchange, and
  ABCG's CG iteration execute and time in each iteration (\cref{sec:codepaths});
\item A characterization of the run-time distributions, including how far they depart from
  normality and why (\cref{sec:distributions});
\item A calibration of normal-theory and rank-based tests and intervals against the
  data's own distribution, measuring false positive rate, power, coverage, sample-size
  predictions, and replication across allocations (\cref{sec:statistics});
\item An evaluation, with those tests, of whether each stand-in reproduces the strategy
  verdicts of the run above it, from the job-wide Vernier replay up to CG
  (\cref{sec:spmvcg,sec:replay,sec:vernier});
\item A set of bounded recommendations for comparing communication strategies on shared
  systems (\cref{sec:discussion}).
\end{itemize}
\Cref{sec:background} describes the stand-ins and the strategies, \cref{sec:setup} the
experiment, \cref{sec:related} related work, and \cref{sec:conclusion} concludes.

\section{Background}
\label{sec:background}

The remainder of this section describes the three halo-exchange strategies, the seven ways
of running the exchange that the campaign compares, and the measurement each run reports.

\paragraph{Strategies:} Each run uses one of three strategies, which we call \emph{arms}.
\textbf{MPI} is plain MPI point-to-point exchange. \textbf{STANDARD} is MPI Advance's
neighbor alltoallv without aggregation~\cite{mpiadvance2023}. \textbf{LOCALITY} is MPI
Advance's locality-aware neighbor alltoallv, which gathers each NUMA region's messages to
another region at a leader rank and sends them as one message (\texttt{-s NUMA}); it is
the only aggregating arm.

\paragraph{ABCG CG:} ABCG's \texttt{benchmark\_cg\_pilot} runs conjugate gradient on a
matrix partitioned into one contiguous block of rows per rank~\cite{abcg}. We treat its time
per iteration as the run the stand-ins are meant to predict. On all six matrices CG
reached ABCG's 500-iteration cap without converging, so its times are per fixed iteration,
not time to solution.

\paragraph{ABCG SpMV exchange:} ABCG's \texttt{benchmark\_spmv} runs only CG's
pack-and-exchange code, once per iteration, followed by a barrier.

\paragraph{\pmrep:} The irregular-benchmarking replay benchmark reads the matrix in
PETSc-derived \texttt{.pm} format, builds the same block partition and halo, and exchanges
it through Cabana's halo gather, with a barrier after every exchange and no computation. In
all 2,880 paired runs, the minimum, maximum, and average over ranks of its messages and
bytes sent equal ABCG's exactly.
% src: report.tex sec. 2; results/runs.csv pm vs abcgspmv numberOfmessages/data_sent min/max/avg, 0 mismatches in 2,880 pairs

\paragraph{Fine and coarse replays:} The same benchmark can instead read a Vernier capture
of one ABCG SpMV exchange~\cite{vernier2025}. The \emph{fine} replay reads Vernier's
per-rank hierarchical capture and sends each rank's recorded number of elements to its
recorded partners; its per-rank message and byte counts equal the \pmrep's in every run.
The \emph{coarse} replay reads Vernier's job-wide aggregate histograms and samples each
rank's neighbor count, partners, and message sizes from them, drawing a new pattern in
every run.
% src: audit.csv fine_vs_pm_mismatched_values = 0 at every node count

\paragraph{Non-persistent ABCG:} The campaign also ran ABCG's SpMV exchange and CG with one
non-persistent MPI Advance call per exchange instead of a persistent request. We include
them in the statistical calibration, where they add 144 comparisons, and leave them out of
the stand-in evaluation, as the companion report did, because their LOCALITY exchange takes
0.31 to 54 times as long as the persistent one~\cite{ladderreport}.
% src: ratios.csv stand_in=abcgspmv reference=abcgspmvnp arm=LOCALITY slow_runs=kept (1/max 0.31, 1/min 53.7); 144 = 2 kinds x 72

\paragraph{Metric:} The time per iteration $q$ of a run is its total timed time, maximized
over ranks, divided by its number of timed iterations $n$. Because each timed iteration
ends at a barrier or an allreduce, the minimum and maximum over ranks differ by at most
0.74\%.
% src: audit.csv rank_spread_timedTotal_max max over node counts = 0.007373

\section{What Each Program Runs and Times}
\label{sec:codepaths}

The stand-ins send the same messages as ABCG, but they do not execute the same code. This
section lists what each program does inside one timed iteration, because those differences
are the mechanisms behind the disagreements in \cref{sec:spmvcg,sec:replay}.
\Cref{tab:codepaths} summarizes them.

\begin{table}[t]
\centering
\caption{What one timed iteration contains in each program, at the commits the campaign
ran (irregular-benchmarking 2593e29, ABCG 29953fe). ``Exchange'' is one halo exchange of
the SpMV's ghost values.}
\label{tab:codepaths}
\small
\begin{tabular}{@{}>{\raggedright\arraybackslash}p{0.16\linewidth}>{\raggedright\arraybackslash}p{0.26\linewidth}>{\raggedright\arraybackslash}p{0.26\linewidth}>{\raggedright\arraybackslash}p{0.26\linewidth}@{}}
\toprule
& \pmrep{} (Cabana) & ABCG SpMV & ABCG CG \\
\midrule
Timer & \texttt{high\_resolution\_clock} (system clock) & \texttt{MPI\_Wtime} & \texttt{MPI\_Wtime}, end of one iteration to end of the next \\
Exchanges & 1 & 1 & 1, plus 1 more every 8th iteration \\
Pack & Cabana gather, Kokkos kernel into a send buffer & serial loop into a \texttt{std::vector<double>} & as ABCG SpMV \\
Unpack & scatter into ghost slots & none & none; the off-process SpMV reads the receive buffer \\
Overlap & none & none & on-process SpMV between start and wait \\
Barrier & \texttt{MPI\_Barrier} after each exchange (\texttt{-b}) & \texttt{MPI\_Barrier} after each exchange & none \\
Other & none & none & 2 allreduces, vector updates \\
Persistent setup & first exchange, in the untimed startup window & \texttt{init()}, before the loop & as ABCG SpMV \\
\bottomrule
\end{tabular}
% src: irregular-benchmarking 2593e29 src/benchmark.cpp (clock_type, runSample loop: apply, -b barrier) and scripts/gatherdata.sbatch (-x EXPORT -b -W); ABCG 29953fe benchmarks/benchmark_spmv.cpp (timed loop), benchmarks/benchmark_cg_pilot.cpp (timed_solve, recompute_r = 8), src/cg.cpp (spmv: pack, start, on_proc, wait, off_proc with get_recvbuf())
\end{table}

\paragraph{ABCG CG:} Each CG iteration packs the ghost values, starts the exchange,
multiplies the on-process block of the matrix, waits for the exchange, and multiplies the
off-process block directly from the receive buffer, so CG never copies received values.
It then computes two inner products, each an allreduce, and updates three vectors. Every
eighth iteration (iterations 0, 8, 16, \ldots, 63 of every 500) recomputes the residual
with a second SpMV and therefore a second exchange, so a CG iteration contains 1.126
exchanges on average. No barrier separates iterations, and each iteration is timed from
the end of the previous one. CG is the only one of the three programs that overlaps
communication with computation.
% src: ABCG 29953fe benchmark_cg_pilot.cpp timed_solve; 63 = |{0,8,...,496}|

\paragraph{ABCG SpMV exchange:} Each timed iteration packs the ghost values with a serial
loop, starts and completes the exchange, and enters an \texttt{MPI\_Barrier}. There is no
computation to overlap and no unpack. For STANDARD and LOCALITY, ABCG creates its persistent
MPI Advance request, including the locality-aware communicator and its aggregation plan,
before the timing loop, and passes the global column index of every ghost value.
% src: ABCG 29953fe benchmark_spmv.cpp; src/comm/spmv_comm_mpil_pers.hpp (init: MPIL_Neighbor_alltoallv_init_ext_topo with global indices)

\paragraph{\pmrep:} Each timed iteration calls Cabana's \texttt{gather.apply()} on a
one-double-per-entry AoSoA and then the benchmark's \texttt{MPI\_Barrier}. The MPI arm uses
Cabana's plain MPI halo and STANDARD and LOCALITY use Cabana's \texttt{LocalityAware}
halo, which builds MPI Advance's persistent neighbor alltoallv on the first exchange.
Because the benchmark times 25 startup exchanges apart from the $n$ timed ones, that setup
falls outside $q$.
% src: irregular-benchmarking 2593e29 src/benchmark.cpp runSample (window vs timed totals) and main (Cabana::Mpi vs Cabana::LocalityAware); run headers "-alltoallv: NEIGHBOR_ALLTOALLV_INIT_*"

\section{Experimental Setup}
\label{sec:setup}

\paragraph{Machine and software:} We ran on Dane at LLNL, whose nodes have two Intel Xeon
Platinum 8480+ sockets and two Omni-Path host fabric interfaces, with 112 MPI ranks per
node, each bound to its own core, at 2, 4, 8, and 16 nodes (224 to 1,792 ranks), in the
shared \texttt{pbatch} queue. We excluded nodes dane261 and dane764, which report 2 NUMA
domains instead of 8 and crash LOCALITY. The software was Open MPI 4.1.2, gcc 13.3.1, MPI
Advance \texttt{locality\_aware} commit 752a19b, irregular-benchmarking 2593e29, ABCG
29953fe, and Vernier 09d650b.
% src: report.tex sec. 3; compute-node facts from the 2026-10-15 write test (10-15-2026.txt: Omni-Path hfi1_0/hfi1_1, Xeon Platinum 8480+)

\paragraph{Matrices:} We used six SuiteSparse matrices: crankseg\_2 (63,838 rows, 14.1M
nonzeros), nd12k (36,000 rows, 14.2M), x104 (108,384 rows, 10.2M), crankseg\_1 (52,804
rows, 10.6M), m\_t1 (97,578 rows, 9.8M), and nd24k (72,000 rows, 28.7M). We chose them
for the structurally different halo exchanges their SpMV produces; each CG iteration
repeats the same exchange, so convergence is not needed for timing.
% src: report.tex sec. 3 (run file headers); rationale: 10-15-2026.txt sec. 6 item 9

\paragraph{Procedure:} For each node count, a pilot job captured each matrix's SpMV halo
with Vernier and then ran every kind on every matrix on STANDARD to choose an iteration
count $n$ that takes about 10\,s; all three arms use that $n$. Each run warms up
connections to ranks $r\pm1$ without timing, performs $k=25$ startup iterations, which $q$
excludes, and then $n$ timed iterations. Forty launches per node count, each a separate
Slurm allocation, ran every combination of 7 kinds, 6 matrices, and 3 arms, 126 runs per
launch, each launch in its own random order with one \texttt{srun} per run. At most four
launches of one node count ran at once, and up to 16 overall; the 160 launches ran over
23.7 hours. Every run exited normally, every benchmark MPI Advance run verified its ghost
values against plain MPI, and no ABCG SpMV run reported a wrong halo value.
% src: report.tex sec. 3; audit.csv

\paragraph{Comparisons:} The unit of replication is the launch, because runs within a
launch share an allocation. A \emph{comparison} is a pair of arms, $A$ and $B$, of one kind
on one matrix at one node count; its data are the 40 paired values
$d = \ln(q_A/q_B)$, one per launch, and $d<0$ means $A$ was faster in that launch. There are
504 comparisons (4 node counts, 7 kinds, 6 matrices, 3 arm pairs). We say a run
\emph{stalled} if its time exceeds the fastest launch's of its configuration by at least
one whole 0.160\,s step in a configuration where those steps are present
(\cref{sec:distributions}), or if its $q$ exceeds five times its configuration's median.
We classify each comparison by its \emph{stall regime}: \emph{none} if neither arm has a
stalled run (321 comparisons), \emph{minority} if at most half of either arm's launches
stalled (101), and \emph{majority} otherwise (82).
% src: estimands.csv regime counts

\section{How the Run Times Are Distributed}
\label{sec:distributions}

\begin{figure}[t]
\centering
\includegraphics[width=0.75\linewidth]{stall_quantum_excess_fine_crankseg_2_STANDARD_16nodes.png}
\caption{Extra timed time of each of 40 launches of the fine replay of crankseg\_2 on
STANDARD at 16 nodes ($n=57$ timed iterations), relative to the fastest launch, in units of
0.160\,s. Every bar falls on a whole number.}
\label{fig:stallhist}
\end{figure}

\begin{figure}[t]
\centering
\includegraphics[width=0.7\linewidth]{stall_steps_pm_crankseg_2_4nodes.png}
\caption{Extra timed time of each launch of the \pmrep{} of crankseg\_2 at 4 nodes, per arm,
in units of 0.160\,s ($n=440$ on every arm). MPI and STANDARD lose up to 432 stalls in 440
iterations; LOCALITY loses none.}
\label{fig:stallsteps}
\end{figure}

Shapiro-Wilk rejects normality at the 5\% level in 86\% of the 504 configurations of node
count, kind, matrix, and arm on $q$, in 62\% on $\ln q$, and in 61\% of the 504 paired
log ratios $d$. The median skewness is 1.24 on $q$ and 0.70 on $\ln q$. A diagnostic test
alone says little at 40 observations per configuration, because Shapiro-Wilk detects
departures too small to matter; the mechanism behind the departure is what decides which
statistics fail.
% src: normality.csv (cell q 0.859, cell ln q 0.621, comparison d 0.607; median skew 1.240, 0.695)

The mechanism is a timing anomaly that the companion report identified: some MPI and
STANDARD runs lose whole multiples of 0.160\,s inside the timed loop. In 61 of the 447
configurations whose total timed time varies by at least 0.3\,s across launches, the
launch-to-launch differences line up on a period with a median of 0.1600\,s and phase
coherence up to 0.995 (\cref{fig:stallhist}). Among the replay and persistent ABCG kinds,
266 of 9,600 MPI and STANDARD runs exceed five times their configuration's median, and no
LOCALITY run does. Where stalls are common they are not outliers but most of the time. At
4-node crankseg\_2, the \pmrep's MPI and STANDARD runs lose up to 432 stalls in 440
iterations, and its LOCALITY runs none (\cref{fig:stallsteps}). Neither launch start time,
any single node, nor the preceding run predicts which launches stall, and the cause is
unknown~\cite{ladderreport}.
% src: stall_quantum.csv (447 rows, 61 > 0.5, best period median 0.1600, max 0.9952); slow_runs.csv / runs.csv (266 of 9,600; 0 LOCALITY); stall_steps figure (pm 4n crankseg_2, max 432, n = 440)

A stall distribution of this kind is a mixture: a narrow fast component, plus a component
shifted by $k\times0.160$\,s whose weight differs between arms and launches. A mixture
separates the estimands that different tests target. The median of $d$ describes the
typical launch, the mean of $d$ (what a paired $t$-test on $\ln q$ targets) weights the
stalled launches by their log time, and the ratio of mean times, $\ln(\bar q_A/\bar q_B)$,
is the expected cost including stalls. These three agree in sign in 448 of the 504
comparisons. Where they disagree, the estimates sit near zero, and in no comparison do the
95\% intervals of the median and of the mean, or of the median and of the ratio of means,
exclude zero on opposite sides (\cref{fig:estimands}). Removing whole stalls, however, changes the answer
where most launches stall: in 23 of the 82 majority-regime comparisons the faster arm by
median changes once each run's whole stalls beyond the fastest launch's are subtracted. At
4-node crankseg\_2, the \pmrep's median STANDARD time is 16.6 times its LOCALITY time; with
stalls removed it is 0.74 times. Because the fastest launch may itself contain stalls, this
correction is a lower bound on the stall time.
% src: estimands.csv (sign(median_lr) = sign(mean_lr) = sign(rom_lr) in 448; confident flips 0; clean_changes 23 of 82 majority; pm 4n crankseg_2 STANDARD/LOCALITY median 16.591, clean 0.741)

\begin{figure}[t]
\centering
\includegraphics[width=0.62\linewidth]{estimand_median_vs_rom.png}
\caption{Median over 40 launches of $d=\ln(q_A/q_B)$ against the log ratio of the arms' mean
times, for all 504 comparisons, colored by stall regime. Points in the upper-left and
lower-right quadrants pick different arms; all of them lie near the origin.}
\label{fig:estimands}
\end{figure}

\section{What Normal-Theory Statistics Get Right and Wrong}
\label{sec:statistics}

A test is reliable for this purpose if (1) it rejects a true null hypothesis no more often
than its nominal level, (2) it detects a real difference of a size that matters, (3) its
intervals contain the quantity they estimate at their nominal rate, and (4) its verdicts
hold on a second set of allocations. We measured all four for two normal-theory tests, the
paired $t$-test on $\ln q$ and on $q$, and two rank-based tests, the exact sign test on $d$
and the Wilcoxon signed-rank test~\cite{wilcoxon1945individual}, whose estimand is the
Hodges-Lehmann pseudo-median~\cite{hodges1963estimates}.

\paragraph{Method:} To judge the tests without making any one of them the referee, we use
each comparison's own 40 values as the population. For each test we shift the values so
that the test's own estimand is exactly zero (the mean of $d$ for the $t$-test on $\ln q$,
the mean of $q_A-q_B$ for the $t$-test on $q$, the median for the sign test, the
pseudo-median for Wilcoxon), draw $n=5$, 10, 20, or 40 launches with replacement 2,000
times, and count rejections at the 5\% level. Sampling with replacement from the
comparison's values keeps the shape of its distribution, including its stalls, and avoids
any overlap between the samples and the target. The Monte Carlo standard error of a 5\%
rate is about 0.5 percentage points. Power uses the same shifted populations moved by
$\ln 1.05$ or $\ln 1.10$. Coverage uses the unshifted populations. Replication uses the
real launches split into two halves of 20.

\begin{figure}[t]
\centering
\includegraphics[width=0.75\linewidth]{calibration_size_minority.png}
\caption{False positive rate when each test's own null hypothesis is true, for the 101
comparisons with stalls in under half the launches, by number of launches. Points are
medians over comparisons and bars the 10th to 90th percentiles. The grey series is the
$t$-test on $\ln q$ in populations where the typical launch is a tie.}
\label{fig:size}
\end{figure}

\paragraph{False positives:} The paired $t$-test on $\ln q$ keeps close to its nominal
level. Its median false positive rate over comparisons is 5.1 to 6.2\% at every number of
launches and stall regime, and the share of comparisons above 7.5\% falls from 26 to 28\%
at 5 launches to 1 to 8\% at 40. The $t$-test on raw $q$ does not. With stalls present its
median rate is 10 to 14\% at 5 launches and 6 to 7\% at 40, and it exceeds 7.5\% in 63 to
70\% of stalled comparisons at 5 launches and in 21 to 41\% at 40 (\cref{fig:size}). The
sign test never exceeds its nominal level, because a binomial count is exact whatever the
distribution, and Wilcoxon stays at 5.0 to 5.2\% from 10 launches on even though stalls
make $d$ asymmetric.
% src: calibration_size.csv (median by regime x n: t_ln 0.051-0.062; >7.5%: 0.26-0.28 at n=5, 0.01-0.08 at n=40; t_q majority/minority n=5 0.104/0.136, n=40 0.062/0.068; >7.5% 0.63/0.70 at n=5, 0.21/0.41 at n=40; sign <= 0.042; wilcoxon 0.050-0.052 for n >= 10)

\paragraph{The estimand matters more than the distribution:} The grey series in
\cref{fig:size} runs the $t$-test on $\ln q$ in populations where the median of $d$ is
zero, so that the typical launch is a tie. At 40 launches the $t$-test then declares a
winner in a median of 10 to 11\% of draws, and in more than 7.5\% of draws for 59 to 62\%
of comparisons. This is not a false positive for the $t$-test, whose null concerns the
mean; it is how often the $t$-test answers a different question from the one an engineer
choosing an arm for a typical run is asking.
% src: calibration_size.csv t_ln_vs_median n=40 median 0.098-0.110; >7.5% share 0.59-0.62

\begin{figure}[t]
\centering
\includegraphics[width=0.75\linewidth]{calibration_power_10pct.png}
\caption{Probability of detecting a true 10\% difference, by number of launches, test, and
stall regime (median over comparisons).}
\label{fig:power}
\end{figure}

\paragraph{Power:} Without stalls, every test detects a 10\% difference with at least 0.92
probability from 10 launches on, and the $t$-test is the only one that can reject at 5
launches, where the smallest $p$-value of the sign and Wilcoxon tests is 0.0625. With
stalls, no test detects a 10\% difference with more than 0.22 probability even at 40
launches (\cref{fig:power}). Pooled over all comparisons, Wilcoxon and the sign test detect
a 10\% difference with 0.89 and 0.80 probability at 20 launches against 0.72 for the
$t$-test.
% src: calibration_power.csv (10pct: none n=10 t 0.980, sign 0.919, wilcoxon 0.984; n=5 t 0.736; minority/majority n=40 <= 0.220; pooled n=20 wilcoxon 0.887, sign 0.801, t 0.722)

\paragraph{Coverage:} The $t$ interval for the mean of $d$ covers 94.0 to 94.6\% in the
median comparison, but 89 to 93\% in the worst tenth. The distribution-free interval for
the median formed by order statistics covers at least its nominal level by construction
(95.8 to 97.8\% for 10 to 40 launches), except at 5 launches, where no such interval can
reach 95\%. The BCa bootstrap interval~\cite{efron1987better} for the median, which the
companion report and our earlier analyses used, covers 93.0, 93.9, and 94.4\% at 10, 20,
and 40 launches in the median comparison, slightly below nominal.
% src: calibration_coverage.csv (t median 0.940-0.946, p10 0.893-0.930; order statistic 0.958-0.978 for n=10..40, 0.938 at n=5; BCa 0.930/0.939/0.944, 60 comparisons)

\paragraph{Sample size:} We chose 40 launches with a normal-theory calculation of how many
launches pin the mean of $\ln q$ to $\pm10\%$ at 95\%. The same calculation asks for more
than 40 launches in 187 of the 504 configurations, and up to 1,111. At the number of
launches it asks for, the $t$ interval covers less than 93\% in 172 configurations and as
little as 77\% (\cref{fig:samplesize}). For the median, the distribution-free interval
needs more launches than the normal-theory count in 215 of the 267 configurations where
that count is at least 6, and more than 400 in 95 of them; part of that difference is the
median's lower efficiency, not a failure of the calculation.
% src: sample_size.csv (n_normal > 40: 24+82+81 = 187; max 1,111; t coverage < 0.93 in 172; min 0.77; n_median > n_normal in 215 of 267, > 400 in 95)

\begin{figure}[t]
\centering
\includegraphics[width=0.62\linewidth]{sample_size_normal_vs_median.png}
\caption{Launches needed for a $\pm10\%$ interval, by a normal-theory calculation for the
mean of $\ln q$ (horizontal) and for the distribution-free interval of the median
(vertical), per configuration. Grey lines mark the campaign's 40 launches.}
\label{fig:samplesize}
\end{figure}

\paragraph{Replication:} The four tests replicate their own significant verdicts on the
other half of the launches equally well, 92.3 to 93.3\% of the time, and none reverses a
verdict when the first 20 launches to start are compared with the last 20
(\cref{fig:replication}); over 500 random splits, 0.03\% of significant verdicts reverse.
With stalls in under half the launches, the sign test reaches significance less often than
the $t$-test (35\% of half-samples against 45\%) and replicates less often (73 to 76\%
against 85 to 89\%). Paired $d$ shows no trend with launch start time, position in the
launch, or the number of concurrently running launches in any comparison (Spearman
correlation, Benjamini-Hochberg~\cite{benjamini1995controlling} $q<0.05$), and a bootstrap
over blocks of five consecutive launches gives intervals of the same width as the ordinary
launch bootstrap (median ratio 1.01), so the launches behave as exchangeable for paired
comparisons.
% src: replication.csv summary (replication 0.9231-0.9333; chronological reversals 0; random 0.0003; minority sign sig_rate 0.3431-0.3515 vs t_ln 0.4455-0.4460; replication 0.7324-0.7583 vs 0.8538-0.8889); exchangeability.csv (0 of 504 at q < 0.05 for each covariate; block_over_iid median 1.01)

\begin{figure}[t]
\centering
\includegraphics[width=0.72\linewidth]{replication_chronological.png}
\caption{For each test, the share of verdicts on 20 launches that are significant, and of
those, the share the other 20 launches confirm or reverse, splitting the 40 launches into
the first and last 20 to start.}
\label{fig:replication}
\end{figure}

\paragraph{Verdicts on all 40 launches:} On the full data, the $t$-test on $\ln q$ and the
sign test, each corrected for multiple comparisons within a kind, agree on 472 of the 504
comparisons and never reach opposite verdicts. Of the 32 disagreements, 18 are in
comparisons with stalls, where the $t$-test rejects alone in 14, and 14 are in comparisons
without stalls, where the sign test rejects alone in 12. The two tests name a different
best arm in 16 of the 168 configurations of kind, matrix, and node count. The $t$-test on
raw $q$ disagrees with the sign test on 46 comparisons.
% src: verdicts.csv crosstab (166+126+180 = 472); verdict_disagreements.csv (stalls: t only 14, sign only 4; power: sign only 12, t only 2); best_arm_tests.csv (16 of 168); t_q crosstab off-diagonal 46

\section{Whether ABCG's SpMV Exchange Predicts CG}
\label{sec:spmvcg}

We say a stand-in \emph{reaches CG's verdict} on an arm pair at one matrix and node count
if the sign test, run separately on each, gives the same outcome: $A$ faster, $B$ faster, or
no decision. It \emph{reproduces CG's ratio} if the per-launch difference between the two
kinds' log arm ratios has a 90\% order-statistic interval of its median inside $\pm10\%$,
which is an equivalence test at the 5\% level.

\begin{figure}[t]
\centering
\includegraphics[width=0.9\linewidth]{tracking_summary.png}
\caption{For each stand-in and the run it is compared with, the share of the 72
combinations of node count, matrix, and arm pair at which the stand-in reproduces the arm
ratio within $\pm10\%$ (green), shifts it significantly (red), or reaches the same sign-test
verdict (blue).}
\label{fig:tracking}
\end{figure}

ABCG's SpMV exchange reaches CG's verdict on 62 of the 72 arm pairs (\cref{fig:tracking}).
The companion report found that choosing CG's arm from the SpMV exchange's medians costs CG
at most 4.2\%~\cite{ladderreport}. However, the SpMV exchange reproduces CG's ratio on only
25 of 72 pairs, and only 7 of 24 for each pair involving LOCALITY, where the SpMV
exchange's arm ratio is 0.58 to 1.41 times CG's. The exchange alone therefore orders the
arms correctly but does not predict by how much one beats another.
% src: tracking.csv stand_in=abcgspmv reference=abcgcg (verdict_agree 62; tracks 25; MPI/LOCALITY 7, STANDARD/LOCALITY 7; median_e range 0.578-1.406); report.tex tab:ladder

Eight of the ten disagreements are between MPI and STANDARD at 2 and 4 nodes, on nd12k,
x104, m\_t1, and nd24k. At four of them (2-node nd12k and x104, 4-node x104 and m\_t1) the
two reach opposite significant verdicts: the SpMV exchange finds STANDARD faster and CG
finds MPI faster, and the two kinds order MPI and STANDARD the same way in only 12.5 to
42.5\% of launches. These are among the eight configurations where the companion report
found every stand-in preferring STANDARD while CG is fastest on MPI, and none of them
contains a stalled run.
% src: tracking.csv abcgspmv vs abcgcg rows with verdict_agree False (8 MPI/STANDARD at 2-4 nodes; opposite at 2n nd12k, x104, 4n x104, m_t1; order_agree 0.125-0.425)

The code paths in \cref{sec:codepaths} give two candidate mechanisms for this bias, which
this campaign cannot separate. First, CG overlaps the exchange with the on-process SpMV.
MPI's nonblocking sends and receives can progress during that multiply, whereas a
neighbor alltoallv implemented as one collective may make no progress until its wait,
so overlap can favor MPI in CG and not in the standalone exchange, which has nothing to
overlap. Second, the SpMV exchange ends every iteration with a barrier and CG does not, so
any difference in how the two arms leave ranks skewed is absorbed by the barrier in one and
carried into the next iteration in the other.

\section{Whether the Replay Benchmark Predicts ABCG}
\label{sec:replay}

\begin{figure}[t]
\centering
\includegraphics[width=0.9\linewidth]{tracking_pm_vs_abcgspmv.png}
\caption{Ratio of the \pmrep's arm ratios to ABCG SpMV's, per node count, matrix, and arm
pair: median over 40 launches with 90\% order-statistic intervals. The grey band is
$\pm10\%$.}
\label{fig:pmspmv}
\end{figure}

The \pmrep{} sends exactly ABCG's messages and bytes, so its disagreements with ABCG come
from the code path. On LOCALITY it is 1.49 to 1.99 times slower than ABCG's SpMV exchange
at every matrix and node count. On MPI and STANDARD it takes 0.95 to 1.14 times ABCG's time
on m\_t1 and x104 at every node count, but 0.95 to 3.79 times on the crankseg matrices,
nd12k, and nd24k, and more than 1.2 times at 23 of their 32 combinations of node count and
arm; the largest ratios are at 8-node crankseg\_2 (3.2 on MPI and 3.8 on STANDARD), and
most of the gap remains when slow runs are dropped. As a result it reproduces ABCG's ratio for
neither arm pair involving LOCALITY at any of the 24 combinations, and shifts it
significantly at 16 of 24 for each (\cref{fig:pmspmv}). It reaches the SpMV exchange's
verdict on 53 of 72 arm pairs and CG's on 49. At 4-node nd12k and 16-node m\_t1 and x104
it finds MPI or STANDARD significantly faster than LOCALITY where ABCG's exchange finds
LOCALITY significantly faster; these are the configurations where the companion report
found that the \pmrep's choice costs CG 1.6 to 21\%.
% src: ratios.csv stand_in=pm reference=abcgspmv slow_runs=kept (LOCALITY 1.488-1.992; MPI/STANDARD m_t1, x104 0.95-1.14; other four matrices 0.950-3.786, > 1.2 at 23 of 32; 8n crankseg_2 3.19/3.79; slow_runs=dropped mostly unchanged, 16n crankseg_2 1.82/2.51 -> 1.35/1.48); tracking.csv pm vs abcgspmv (tracks 0 for LOCALITY pairs, shifted 16 each, verdict_agree 53); pm vs abcgcg verdict_agree 49; report.tex (4n nd12k 21%, 16n m_t1 7.1%, x104 1.6%)

Because the MPI arm shares Cabana's pack, unpack, and barrier with the other arms, those
costs cannot explain a gap that appears only on LOCALITY and is absent from MPI on the
lighter matrices. The gap is therefore in how the benchmark's Cabana path drives MPI
Advance's locality-aware exchange, rather than in Cabana or Kokkos as such.

\section{Whether the Vernier Replays Predict the Matrix Replay}
\label{sec:vernier}

The fine replay reproduces the \pmrep's arm ratio on 38 of 72 arm pairs and shifts it
significantly on only 2, and it reaches the same verdict on 64. Six of its eight
disagreements are between MPI and STANDARD, and the other two are the LOCALITY pairs at
2-node crankseg\_2, a configuration with stalls in under half the launches. The coarse replay reproduces the fine replay's ratio on 16 of 72 pairs and
shifts it significantly on 49, and on the two LOCALITY pairs it reproduces the ratio at
only 1 and 2 of 24 combinations, with ratios from 0.027 to 1.85 times. It reaches the fine
replay's verdict on 52 of 72, and 14 of its 20 disagreements are at 2 and 4 nodes. These
results are consistent with the companion report's finding that the job-wide histograms
lose the partner structure that decides the aggregating arm at small node
counts~\cite{ladderreport}.
% src: tracking.csv fine vs pm (tracks 38, shifted 2, verdict_agree 64; 6 of 8 disagreements MPI/STANDARD, 2 at 2n crankseg_2); coarse vs fine (tracks 16, shifted 49, LOCALITY pairs tracks 1 and 2, median_e 0.027-1.845, verdict_agree 52, 14 of 20 at 2-4 nodes)

\section{Discussion}
\label{sec:discussion}

The calibration and the stand-in comparisons lead us to the following conclusions about
comparing halo-exchange strategies on a shared system like Dane.

\paragraph{1) Test on log time, paired within an allocation:} The paired $t$-test on
$\ln q$ held its false positive rate within about one point of nominal and replicated as
well as the rank-based tests, even though Shapiro-Wilk rejected normality in most
configurations. The $t$-test on raw time did not: with stalls it raised false positives to
10 to 14\% at 5 launches. This recommendation holds for paired comparisons within an
allocation; we did not test comparisons across allocations, such as across node counts,
which do not cancel shared machine state.

\paragraph{2) Choose the estimand before the test:} In stall-dominated configurations the
mean, the median, and the expected cost answer different questions, and removing whole
stalls changed the median's winner in 23 of 82 majority-regime comparisons. A test cannot
make that choice. Where the question is the typical run, the sign test with an
order-statistic interval answers it exactly; where the question is the expected cost of a
long run, the ratio of means with a launch bootstrap does, stalls included.

\paragraph{3) Do not trust normal-theory sample sizes on stall-prone patterns:} The
calculation that chose 40 launches asks for hundreds on the MPI and STANDARD arms of the
stall-prone matrices, and at its own recommended counts its intervals undercover in a
third of configurations. With stalls present, no test detected a 10\% difference reliably
at 40 launches, so the campaign can rank the arms in those configurations only where they
differ by much more.

\paragraph{4) Use the application's own exchange as the stand-in:} ABCG's SpMV exchange
reached CG's verdict on 62 of 72 arm pairs and costs CG at most 4.2\% when it chooses.
However, it does not reproduce CG's arm ratios within 10\% for most LOCALITY pairs, and it
consistently prefers STANDARD where CG is fastest on MPI.

\paragraph{5) A benchmark's code path can matter more than its pattern:} The \pmrep{} sends
the exact pattern and still disagrees with ABCG wherever LOCALITY's advantage is small,
because its locality-aware path is up to twice as slow. A replay benchmark is useful for
choosing whether to aggregate only when the aggregating arm's advantage is larger than the
difference between the benchmark's and the application's implementation of that arm.

\section{Related Work}
\label{sec:related}

\paragraph{Statistics for performance measurement:} Hoefler and Belli's rules for
reporting parallel performance include not assuming normality without diagnostic checking,
and note that nonparametric metrics such as the median assume no particular
distribution~\cite{hoefler2015scientific}. Hunold and Carpen-Amarie showed that MPI
benchmark results vary between \texttt{mpirun} invocations and recommended treating the
invocation as a factor, with 10 to 30
invocations~\cite{hunold2016reproducible,hunold2015mpibenchmarking}. Georges et al.\ and
Kalibera and Jones developed rigorous designs for managed-runtime
benchmarks~\cite{georges2007statistically,kalibera2013rigorous}. Maricq et al.\ measured
variability on testbed hardware and advise nonparametric confidence
intervals~\cite{maricq2018taming}, Duplyakin et al.\ found performance regimes that evolve
over time in years of cloud benchmark data~\cite{duplyakin2020datacenter}, and Agarwal et
al.\ recommend bootstrap intervals when only a few runs are available~\cite{agarwal2021deep}.
Our work is complementary to these: rather than recommending a method, it calibrates
normal-theory and rank-based tests against the measured distribution of one HPC campaign,
and finds that the paired $t$-test on log time survives a departure from normality that
diagnostic tests flag in most configurations, while the $t$-test on raw time and
normal-theory sample sizes do not.

\paragraph{Variability on shared HPC systems:} Skinner and Kramer and Bhatele et al.\
attributed run-time variability on HPC systems to contention from other jobs and system
activity~\cite{skinner2005understanding,bhatele2013there}. The 0.160\,s stalls here are
not predicted by the time, the nodes, or the preceding run, and we have not identified
their source.

\paragraph{Benchmarks as stand-ins:} Collom et al.\ compared locality-aware and standard
neighbor collectives on SpMV and AMG patterns~\cite{collom2023}, and Vernier was designed to
capture application patterns for replay at scale~\cite{vernier2025}. Unlike those studies,
this work scores each stand-in against the run above it with paired tests and traces its
disagreements to differences in executed code.

\section{Conclusions and Future Work}
\label{sec:conclusion}

In this paper we evaluated, on 40 allocations at each of four node counts of Dane, which
statistical tests give reliable verdicts on halo-exchange strategy comparisons, and how
well each of five stand-ins reproduces the verdicts of the run above it. Normality was
rejected in 86\% of configurations on run time, because MPI and STANDARD runs lose whole
multiples of 0.160\,s. Nevertheless, the paired $t$-test on log time kept its false
positive rate at 5.1 to 6.2\% in the median comparison and replicated its verdicts as well
as the sign and Wilcoxon tests. However, the $t$-test on raw time reached 10 to 14\%
false positives at 5 launches with stalls present, normal-theory sample sizes undercovered
in a third of configurations, and no test detected a true 10\% difference reliably in
configurations with stalls. ABCG's SpMV exchange reached CG's verdict on 62 of 72 arm pairs
but reproduced CG's arm ratios within 10\% on only 25, and it preferred STANDARD where CG
was fastest on MPI. The \pmrep{} reached CG's verdict on 49 of 72 arm pairs, because its
locality-aware path was 1.49 to 1.99 times slower than ABCG's on the same messages.

Several limitations of this study suggest future work. First, every run reports only
totals over its iterations, so the stall mixture is inferred from launch-to-launch
differences rather than seen in single iterations. Second, the campaign covers one machine,
one MPI library, and allocations spanning one day. Third, we identified candidate
mechanisms for the MPI versus STANDARD bias and for the replay's slower locality-aware path
from the code but did not isolate them, which requires runs that remove one difference at
a time, such as an unpack added to ABCG or a replay without Cabana.

\bibliographystyle{plain}
\bibliography{references}

\end{document}
